% Manuscript v5: targeted corrections following referee round 2.
\documentclass[11pt]{article}
%TC:envir table 0 1
%TC:envir tabular 1 1
%TC:envir longtable 1 1
\usepackage[T1]{fontenc}
\usepackage[utf8]{inputenc}
\usepackage{lmodern}
\usepackage[margin=25mm]{geometry}
\usepackage{amsmath,amssymb,graphicx,booktabs,longtable,array}
\usepackage[authoryear,round]{natbib}
\usepackage[hidelinks]{hyperref}
\usepackage{url,placeins}
\setlength{\emergencystretch}{2em}
\newcommand{\dd}{\mathrm d}
\DeclareMathOperator{\Var}{Var}
\DeclareMathOperator{\erf}{erf}
\DeclareMathOperator{\erfc}{erfc}
\DeclareMathOperator{\erfi}{erfi}
\DeclareMathOperator{\erfcx}{erfcx}

%TC:ignore
\title{Quantum and classical physics-informed networks for laser-pulse
thermoelasticity: nominal and output-affecting parameter comparisons}
\author{Hussain Al-Qahtani\thanks{Corresponding author:
\href{mailto:ahtanih@kfupm.edu.sa}{ahtanih@kfupm.edu.sa}.}\\
\small Department of Mechanical Engineering,\\
\small King Fahd University of Petroleum and Minerals,\\
\small Dhahran 31261, Saudi Arabia}
% Further co-authors may be added here after their contributions are confirmed.
\date{}
%TC:endignore
\begin{document}
%TC:ignore
\maketitle
%TC:endignore

\begin{abstract}
Laser-pulse heating couples thermal diffusion to elastic-wave propagation,
providing a demanding test of physics-informed learning. We compare quantum
and classical networks against corrected exact thermoelastic fields, using
internally prespecified research questions and protocol, with execution
changes disclosed. Under the executed recipe, the classical multilayer
perceptron has lower mean temperature and stress errors in every nominal-count
and output-affecting-count comparison (Tables~\ref{tab:e1} and \ref{tab:e1b}).
The larger circuit takes 171.7--228.3 times the nominally matched network's
recorded CPU runtime (Table~\ref{tab:e1}). The full physics objective raises
one selected gradient variance, but finite-range and tail slopes establish
neither asymptotic decay nor barren-plateau mitigation
(Table~\ref{tab:e2summary}; Supplementary Table~S9).
Across the tested parameter and collocation budgets and the production
training-label sweep conditional on separately tuned hyperparameters,
the circuit shows no stress-efficiency advantage over the multilayer perceptron
(Table~\ref{tab:e3}). Supervision lowers the circuit's attainable stress
error but leaves a classical accuracy gap (Table~\ref{tab:e3}).
Secondary wave speeds preserve the circuit--perceptron ordering
(Table~\ref{tab:gamma}). Common-grid inference controls show greater stress
degradation at the lowest shot budget than under either tested
exact-expectation depolarizing channel (Table~\ref{tab:e4ratios}); trained
readout gains amplify sampling fluctuations. The shared initial-rate
exclusion leaves part of the continuum stress evolution to model-dependent
implicit selection. Together with limited replication of the efficiency
sweeps and contention-affected timings, this restricts the comparison to
the declared recipe. The study establishes no quantum advantage.
\end{abstract}

\section*{Highlights}
\begin{itemize}
\item Nominal and output-affecting count comparisons favor the classical MLP.
\item No tested production-budget stress comparison favors the circuit.
\item Finite-range gradient slopes do not establish asymptotic decay.
\item Trained readout gains amplify finite-shot inference errors.
\item Initial-rate exclusions leave part of the stress evolution unconstrained.
\end{itemize}
\noindent\textbf{Keywords:} physics-informed neural networks; quantum machine
learning; laser heating; thermoelasticity; gradient variance; readout noise.

\section{Introduction}

Laser absorption produces localized heating, constrained thermal expansion,
and propagating stress. Predicting these coupled fields supports pulse
selection for laser processing, thermal-stress control, and interpretation
of thermal barrier coating tests. The semi-infinite pulse-heating model of
\citet{AlQahtani2026} offers an exact-reference benchmark for this purpose.
Its simplicity is deliberate: a homogeneous half-space cannot predict
coating-interface failure, plasticity, or material removal. More elaborate
thermoelastic descriptions include non-Fourier transport and strain
feedback \citep{WangXu2001}; the present benchmark isolates Fourier
conduction and one-way thermal forcing of a stress wave.

Physics-informed neural networks (PINNs) approximate fields while penalizing
equation and boundary residuals \citep{Raissi2019}. Thermal-mechanical
learning for laser metal deposition provides an application precedent
\citep{Sharma2024}. Accurate stresses, rather than an aggregate training
loss alone, matter in such applications. Exact fields permit separate
assessment of temperature, stress, and computational expense without
uncertainty from a learned reference.

Differentiable quantum circuits can solve differential equations
\citep{Kyriienko2021}; repeated data encoding enlarges their function
representation \citep{PerezSalinas2020}. Quantum and hybrid PINNs have been
compared with classical networks \citep{Trahan2024}. In particular,
\citet{Berger2025} report favorable comparisons using trainable classical
embeddings and equal parameter counts. Heat-equation and parabolic-PDE
studies \citep{Dehaghani2025,Tran2026}, adaptive loss weighting
\citep{Setty2025}, and a recent embedding comparison
\citep{TranEmbedding2026} establish substantial prior work. A quantum
physics-informed graph-network proposal for laser cutting also exists
\citep{Ruhi2023}. The contribution here is a controlled thermoelastic
benchmark, not a new general quantum-PINN method or a broad priority claim
for laser heating.

Trainability and prediction accuracy are separate questions. Barren-plateau
results depend on circuit ensembles, depth, cost locality, and noise
\citep{McClean2018,Cerezo2021,WangNoise2021}. A PDE-based mitigation proposal
uses qubit-indexed finite-difference field values \citep{Ukwatta2026}; our
coordinate-encoded circuit instead differentiates fixed observable
readouts. Consequently, its gradient diagnostic does not directly test
that proposal's representation or locality argument. Increasing a loss's
scale also increases gradient variance without establishing improved
trainability.

Classical controls are essential to a meaningful quantum comparison
\citep{Bowles2024}. Encoding-dependent Fourier representations
\citep{Schuld2021}, classical surrogates of trained quantum models
\citep{Schreiber2023}, and sufficient conditions for random-feature
dequantization \citep{Sahebi2025} motivate a spectral comparator. Here it is
an independently trained classical spectral baseline motivated by
dequantization results, with no claim to reproduce the circuit's learned
frequencies or function.

The internal brief asks:
\begin{enumerate}
\item \textbf{RQ1:} Does a quantum PINN match a classical PINN with the same
parameter count on these fields?
\item \textbf{RQ2:} Does the physics PDE-residual loss change barren-plateau
behavior, measured through gradient variance versus qubit count and depth,
compared with a data loss?
\item \textbf{RQ3:} Is the quantum PINN more parameter-efficient or
data-efficient?
\end{enumerate}
The brief supplies no directional statistical hypotheses, equivalence
margin, or decision rule. We therefore answer descriptively, without
retrospective confirmatory hypotheses. For RQ2, the executed statistic uses
the full physics objective, including boundary and initial-rate terms.
For RQ3, the label comparison concerns production training labels
conditional on separately selected hyperparameters, not total exact-label
acquisition (Section~3.3).
The study follows internally prespecified research questions and protocol,
with the execution changes disclosed in Supplementary Section~S1.
Its contributions are:
\begin{itemize}
\item replicated compact-model comparisons on coupled heat--stress fields,
with nominal and output-affecting parameter accounting;
\item corrected closed forms checked by inverse-Laplace evaluation and an
independent finite-difference formulation;
\item characterization of readout-gain amplification in selected trained
thermoelastic circuits, using common-grid clean and noisy controls.
\end{itemize}
The general link between observable coefficients and sampling variance,
including variance regularization, has prior precedent
\citep{Kreplin2024}. The contribution is its measured manifestation here.

\section{Problem statement and exact solution}

\subsection{Physical assumptions and dimensional conversion}

Consider a homogeneous isotropic half-space with constant coefficients,
Beer--Lambert absorption, an insulated traction-free surface, and
uniaxial strain. Let $X,\tau,\Theta,u,\Sigma$ denote dimensional depth,
time, temperature rise, displacement, and tensile-positive axial stress.
The constitutive law and momentum balance are \citep{WangXu2001}
\begin{align}
 \Sigma&=M u_X-b_{\rm th}\Theta,&
 \rho u_{\tau\tau}&=\Sigma_X,\\
 M&=\lambda+2\mu=\rho c_{\rm el}^2,&
 b_{\rm th}&=\frac{E\alpha_T}{1-2\nu}.
\end{align}
Here $\alpha_T$ is the linear expansion coefficient. Eliminating $u$ gives
\begin{equation}
 \Sigma_{XX}-c_{\rm el}^{-2}\Sigma_{\tau\tau}
 =c_2\Theta_{\tau\tau},\qquad
 c_2=\frac{b_{\rm th}}{c_{\rm el}^2},\qquad
 [c_2]=\mathrm{Pa\,s^2\,m^{-2}\,K^{-1}}.
 \label{eq:dimensional}
\end{equation}
The coefficient $b_{\rm th}/(\rho c_{\rm el}^2)=b_{\rm th}/M$, with units
$\mathrm{K^{-1}}$, belongs to the displacement equation multiplying
$\Theta_X$; it is not the coefficient in Eq.~\eqref{eq:dimensional}.

With absorption coefficient $\delta$, conductivity $k$, absorbed intensity
$I_1$, and thermal diffusivity $\alpha_{\rm th}=k/(\rho c_p)$, define
\begin{equation}
 x=\delta X,\quad t=\alpha_{\rm th}\delta^2\tau,\quad
 \Theta_*=\frac{I_1}{k\delta},\quad T=\frac{\Theta}{\Theta_*},\quad
 \gamma=\frac{c_{\rm el}^2}{\alpha_{\rm th}^2\delta^2}.
\end{equation}
The original stress scale is
$\Sigma_0=I_1c_2c_{\rm el}^2/(k\delta)=b_{\rm th}\Theta_*$.
The benchmark uses $\sigma=\gamma\Sigma/\Sigma_0$, giving the explicit
conversion
\begin{equation}
 \boxed{\Sigma\,[\mathrm{Pa}]=
 \frac{E\alpha_T I_1}{(1-2\nu)\gamma k\delta}\,\sigma .}
 \label{eq:pascal}
\end{equation}
Dimensional material and pulse scales must be supplied before interpreting
a predicted stress in pascals. All fields and coordinates below are
dimensionless; runtimes are in seconds.

The benchmark equations and complete physical initial data are
\begin{align}
 T_t-T_{xx}&=e^{-x}f(t),&
 \sigma_{xx}-\gamma^{-1}\sigma_{tt}&=T_{tt},\label{eq:pdes}\\
 T_x(0,t)&=0,&\sigma(0,t)&=0,\\
 T,\sigma&\longrightarrow0\quad(x\longrightarrow\infty),&
 T(x,0)&=\sigma(x,0)=0,\\
 T_t(x,0^+)&=e^{-x}f(0^+),&
 \sigma_t(x,0^+)&=-\gamma e^{-x}f(0^+),\quad x>0.
 \label{eq:rates}
\end{align}
The stress rate follows from initially zero material velocity and the
constitutive law. The heat equation has no strain feedback. Phase change,
plasticity, ablation, interfaces, and temperature-dependent properties
are excluded.

\subsection{Pulses, reference fields, and interfaces}

Table~\ref{tab:problem} defines the pulses and domain. Appendix~\ref{app:exact}
gives the corrected closed forms implemented in \path{ground_truth.py}.
Inverse-Laplace evaluation \citep{deHoog1982} and an independently
discretized evolution problem check these fields; Supplementary
Section~S7 provides the executed verification records and their limits.

\begin{table}[htbp]
\centering
\caption{Problem and evaluation settings from the executed benchmark.
$H$ denotes the Heaviside function; the pulses are causal.
Energy is $E_0=\int_0^\infty f(t)\,\dd t$.}
\label{tab:problem}
\begin{tabular}{p{0.22\linewidth}p{0.68\linewidth}}
\toprule
Item & Definition or value \\
\midrule
Exponential & $f(t)=A_e e^{-\beta t}$, $A_e=1$, $\beta=2$, $E_0=0.5$ \\
Rectangular (step) & $f(t)=A_s[H(t)-H(t-\Delta)]$,
$A_s=1$, $\Delta=0.5$, $E_0=0.5$ \\
Smooth & $f(t)=A(e^{-\beta_1t}-e^{-\beta_2t})$,
$A=4/3$, $\beta_1=2$, $\beta_2=8$, $E_0=0.5$ \\
Main material ratio & $\gamma=9$; dimensionless wave speed $\sqrt\gamma=3$ \\
Network domain & $x\in[0,6]$, $t\in[0,3]$ \\
Test grid & $241$ spatial nodes on $[0,6]$ and $240$ time nodes on $(0,3]$ \\
E1/E1b replication & Seeds $0,1,2,3,4$ in each cell; $105$ E1 and $60$ E1b runs \\
E3 & Exponential pulse, $\gamma=9$; seeds $0,1$; $66$ completed runs \\
Secondary wave speeds & Exponential pulse, $\gamma=25,49$;
seeds $0,1,2$; $24$ runs \\
\bottomrule
\end{tabular}
\end{table}

Sudden startup launches a derivative discontinuity along
$x=\sqrt\gamma\,t$. The smooth pulse cancels its leading startup kink.
The rectangular pulse also has a temporal interface. At fixed interior
$x$, away from intersection with the pre-existing characteristic, define
$[g]_\Delta=g(x,\Delta^+)-g(x,\Delta^-)$; then
\begin{equation}
 [T]_\Delta=[\sigma]_\Delta=0,\qquad
 [T_t]_\Delta=-A_s e^{-x},\qquad
 [\sigma_t]_\Delta=\gamma A_s e^{-x}.
 \label{eq:jumps}
\end{equation}
These conditions follow from the heat equation and the distributional
stress equation; $w=\sigma+\gamma T$ and $w_t$ remain continuous.
The interface $t=\Delta$ differs from the delayed surface characteristic
$x=\sqrt\gamma(t-\Delta)$. The executed objective has no explicit
interface penalty; sampled smooth residuals alone do not certify these
jump conditions.

\subsection{Truncation and supplied information}

At the outer boundary in Table~\ref{tab:problem}, exact Dirichlet traces
for both fields restrict the half-space solution without an additional
artificial-boundary approximation. Every physics-trained model receives these traces.
``Physics-only'' therefore means no interior training labels; exact
boundary values and separate exact validation labels are still supplied.
The fully prescribed reference problem and the relaxed training
constraints in Section~3.2 must be distinguished.
The data-only E3 sweep varies the production training-label count $N_d$
conditional on separately selected hyperparameters. Its axis excludes
exact labels used for tuning and validation diagnostics; Section~3.3 and
Table~\ref{tab:training} account for that additional information.

\section{Models and training}

\subsection{Circuit and classical comparators}

Inputs are $u=x/3-1$ and $v=2t/3-1$ for the domain in
Table~\ref{tab:problem}. Starting from $|0\cdots0\rangle$, each circuit
layer applies on each qubit
\begin{equation}
 G_{\ell q}=R_Z(\theta_{3,\ell q})R_Y(\theta_{2,\ell q})
 R_Z(\theta_{1,\ell q})R_X(b_{\ell q}v)R_Y(a_{\ell q}u),
 \label{eq:gate}
\end{equation}
with the rightmost gate first, followed by the CZ ring. From the released
code, the ring has a single edge for the smallest diagnostic circuit;
Table~\ref{tab:models} records the convention and initializations. Outputs are
\begin{equation}
 y_k=\alpha_k\langle Z_k\rangle+\beta_k,\quad k\in\{0,1\},\qquad
 P_{\rm nom}=5n_qL+4.
 \label{eq:readout}
\end{equation}
Training uses exact statevectors and automatic differentiation through
complex arithmetic, including the higher-order graphs needed to
differentiate second-order residuals with respect to parameters.

\begin{table}[htbp]
\centering\small
\caption{Architectures and initialization from the released code and
executed coordinate-perturbation audit. $P_{\rm out}$ counts
output-affecting coordinates, not proven independent functional dimension.
Fixed spectral frequencies are not trainable parameters.}
\label{tab:models}
\begin{tabular}{p{0.23\linewidth}p{0.67\linewidth}}
\toprule
Item & Definition \\
\midrule
Circuit & $(n_q,L)=(4,2),(6,4)$;
$P_{\rm nom}/P_{\rm out}=44/30,124/88$ \\
Initialization & $a,b\sim\mathcal N(0,1)$;
$\theta_j\sim\mathcal U(0,2\pi)$; $\alpha_k=1$, $\beta_k=0$ \\
CZ ring & Edges $(q,q+1\bmod n_q)$ for $n_q>2$;
one edge $(0,1)$ for $n_q=2$ \\
MLP & $2\to h\to h\to2$, tanh;
$P=h^2+6h+2$. E1: $h=4,8$, $P=42,114$;
E1b: $h=3,7$, $P=29,93$ \\
Spectral & $z=[u,v]B$, $B\in\mathbb R^{2\times8}$,
$B_{ij}\sim\mathcal N(0,\pi^2)$ fixed;
$[\cos z,\sin z]\to h\to2$, tanh, no phase;
$P=19h+2$. E1: $h=2,6$, $P=40,116$;
E1b: $h=2,5$, $P=40,97$ \\
Reference MLP & $2\to64\to64\to64\to64\to2$, tanh;
$P=12802$ \\
Classical initialization & PyTorch \texttt{Linear} defaults after setting
the run seed; $P_{\rm out}=P_{\rm nom}$ denotes no structurally removed
coordinates, not a rank measurement \\
Matching & Nearest count; prescribed tolerance $\pm10\%$;
all deviations reported, including exceptions \\
Coordinate audit & Model seed $0$; input seed $12345$;
$64$ uniform inputs in $[-1,1]^2$; additive perturbation $+0.37$
per scalar; inert if maximum output change is below $10^{-12}$ \\
\bottomrule
\end{tabular}
\end{table}

The terminal CZ ring commutes with the readouts; final-layer operations
on unmeasured qubits and terminal measured-qubit $R_Z$ rotations cannot
affect them. Backward causal structure removes further coordinates.
The released \path{functional_count.py} and its executed audit list the
inert coordinates underlying Table~\ref{tab:models}; Supplementary
Section~S2 extends the accounting to every E3 circuit. These counts do not
measure Jacobian rank or rule out dependencies among the remaining
coordinates. The perturbation is an angle increment only for rotation
coordinates, since encoding scales and affine readout parameters are also
tested.

E1 compares registered counts; E1b retunes classical networks toward the
circuit's output-affecting counts. The smaller spectral row repeats E1
and is not an output-count match. Its minimum is imposed by the implemented
width search, not the architecture; the larger spectral row also slightly
exceeds the prescribed tolerance (Table~\ref{tab:e1b}). The reference MLP
is an attainable-accuracy control. Equal counts do not imply equal
capacity, derivative cost, memory, or optimization work.

\subsection{Hard initial fields and physics objective}

Every family uses
\begin{equation}
 T_\vartheta=t y_0(u,v;\vartheta),\qquad
 \sigma_\vartheta=t y_1(u,v;\vartheta).
 \label{eq:ansatz}
\end{equation}
This enforces initial fields, but leaves their initial rates trainable.
Writing $\langle\cdot\rangle_C,\langle\cdot\rangle_B,
\langle\cdot\rangle_I$ for sample means over collocation points, shared
boundary times, and retained initial-rate locations, the released code uses
\begin{align}
 r_h&=T_{\vartheta,t}-T_{\vartheta,xx}-e^{-x}f(t),&
 r_s&=\sigma_{\vartheta,xx}-\gamma^{-1}\sigma_{\vartheta,tt}
       -T_{\vartheta,tt},\\
 \mathcal L_h&=\langle r_h^2\rangle_C,&
 \mathcal L_s&=\langle r_s^2\rangle_C,\\
 \mathcal L_B&=\langle T_{\vartheta,x}(0,t)^2\rangle_B
 +\langle\sigma_\vartheta(0,t)^2\rangle_B\notag\\
 &\quad+\langle[T_\vartheta(6,t)-T^*(6,t)]^2\rangle_B
 +\langle[\sigma_\vartheta(6,t)-\sigma^*(6,t)]^2\rangle_B,\\
 \mathcal L_I&=\langle[y_0(x,0)-e^{-x}f(0^+)]^2\rangle_I
 +\left\langle\left[\frac{y_1(x,0)+\gamma e^{-x}f(0^+)}\gamma
 \right]^2\right\rangle_I,\\
 \mathcal L_P&=\mathcal L_h+\mathcal L_s+\mathcal L_B
 +\lambda_I\mathcal L_I.\label{eq:loss}
\end{align}
In the rate terms, $y_k(x,0)$ abbreviates evaluation at the corresponding
scaled coordinates. A star denotes an exact field. All weights and
reductions are specified in Table~\ref{tab:training}; no adaptive weighting
is used. Adaptive weighting has been studied in quantum differential
equation solvers \citep{Setty2025} and is a distinct future comparison.

\begin{table}[htbp]
\centering
\caption{Sampling, training, and validation settings from the execution
record and released code for the reference physics recipe. E3 changes
point budgets and, for data-only training, the objective and optimizer
schedule (Table~\ref{tab:e3}). Exponential-distribution parameters are
means. Boundary counts denote shared time samples for all boundary terms.}
\label{tab:training}
\begin{tabular}{p{0.27\linewidth}p{0.63\linewidth}}
\toprule
Item & Executed setting \\
\midrule
Adam stage & $2000$ steps; $N_c=512$ resampled collocation points per step;
$128$ boundary times sampled with replacement from $4096$ precomputed
exact far-boundary values \\
Collocation mixture & Half uniform on $[0,6]\times[10^{-3},3]$;
half with $x\sim\mathrm{Exp}(\text{mean }1)$ clipped to $[0,6]$
and $t\sim\mathrm{Exp}(\text{mean }0.7)$ clipped to $[10^{-3},3]$ \\
Initial-rate locations & Collocation $x$ values retained for $x\ge0.3$,
then evaluated at $t=0$ \\
Loss weights & Each PDE and boundary component has weight $1$;
$\lambda_I=1$ in all physics runs, including E2 \\
Adam schedule & Cosine decay to $5\%$ of the initial learning rate \\
L-BFGS stage & Configured limits: $1000$ iterations, $1250$ evaluations;
optimizer step size $1$; strong-Wolfe line search; history $50$;
gradient tolerance $10^{-12}$; change tolerance $10^{-14}$;
fixed $2048$ collocation points and $512$ shared boundary times \\
Rate candidates & $\{10^{-3},3\times10^{-3},10^{-2},3\times10^{-2},10^{-1}\}$ \\
Rate selection & $1000$ Adam steps per candidate, seed $100$;
$2000$ fixed random validation points; score $e_T+e_\sigma$;
never the test grid \\
Selection grouping & Independently for each pulse, family, and configuration \\
E3 data-only tuning & Rates selected once at $N_d=256$: $256$ exact
training samples and $2000$ exact validation samples shared by the
$5$ candidate rates and both families; seed $100$ \\
E3 production validation & Diagnostics only; no parameter updates,
early stopping, or checkpoint selection \\
\bottomrule
\end{tabular}
\end{table}

The stress equation requires an initial rate. For nonzero-start pulses,
its interior right limit conflicts with the differentiated free-surface
condition at the initial corner. Excluding an interval, however, is
stronger than excluding that point. The executed cutoff in
Table~\ref{tab:training} omits $45.12\%$ of the squared nonzero initial-rate
profile on the truncated domain, since
\begin{equation}
 \frac{\int_0^{x_{\min}}e^{-2x}\,\dd x}
 {\int_0^{6}e^{-2x}\,\dd x}
 =\frac{1-e^{-2x_{\min}}}{1-e^{-12}}.
\end{equation}
This is a squared-profile fraction, not mechanical energy.

The retained continuum constraints do not uniquely determine stress.
Choose a smooth nonzero $g$ supported inside $(0,x_{\min})$ and add a
homogeneous wave $v$ satisfying
\begin{equation}
 v_{tt}=\gamma v_{xx},\quad v(x,0)=0,\quad v_t(x,0)=g(x),\quad
 v(0,t)=v(6,t)=0.
\end{equation}
It preserves the PDE, stress boundaries, zero initial stress, and every
retained rate condition. This function-space argument does not assert
multiple exact zero-loss realizations in a finite neural architecture.
It establishes that part of the stress evolution is left to
model-dependent implicit selection. The experiment compares a common
declared recipe against the complete exact reference; it does not certify
solution of the fully prescribed initial-boundary-value problem.

Table~\ref{tab:xmin} completes the single-seed cutoff comparison.
Stress sensitivity is nonmonotone and model-dependent: the declared
cutoff gives the lowest tested circuit and reference-MLP stress errors,
but the highest of the matched MLP's tested stress errors. The circuit
and reference MLP therefore do not establish a generally beneficial
cutoff. Temperature ordering differs. None of these outcomes is a
replicated effect estimate. The separate weight ablation lowers stress
error when the rate term is restored, but not uniformly temperature
error (Supplementary Table~S7). Both ablations change the temperature-rate
and stress-rate penalties together.

\begin{table}[htbp]
\centering\footnotesize
\caption{Near-surface initial-rate exclusion ablation, exponential pulse,
$\gamma=9$, seed $0$, $N_c=512$. Every row is an individual-run outcome
($n=1$), with no standard deviation. Learning rates are $0.1$ for VQC/MLP
and $0.01$ for MLP-ref. The replay verifies the reference-MLP thread
setting and is not an extra cutoff replicate.}
\label{tab:xmin}
\setlength{\tabcolsep}{3pt}
\begin{tabular}{lrrrr}
\toprule
Model & $P_{\rm nom}/P_{\rm out}$ & $x_{\min}$ & $e_T$ & $e_\sigma$ \\
\midrule
VQC & $124/88$ & $0$ & $0.1824$ & $0.9644$ \\
VQC & $124/88$ & $0.1$ & $0.1709$ & $1.1081$ \\
VQC & $124/88$ & $0.3$ & $0.1939$ & $0.8785$ \\
\midrule
MLP & $114/114$ & $0$ & $0.1330$ & $0.2695$ \\
MLP & $114/114$ & $0.1$ & $0.1294$ & $0.2393$ \\
MLP & $114/114$ & $0.3$ & $0.1959$ & $0.3409$ \\
\midrule
MLP-ref & $12802/12802$ & $0$ & $0.0176$ & $0.1001$ \\
MLP-ref & $12802/12802$ & $0.1$ & $0.0143$ & $0.0705$ \\
MLP-ref & $12802/12802$ & $0.3$ & $0.0085$ & $0.0496$ \\
MLP-ref & $12802/12802$ & $0.6$ & $0.0166$ & $0.0617$ \\
MLP-ref & $12802/12802$ & $1$ & $0.0082$ & $0.0732$ \\
MLP-ref replay & $12802/12802$ & $0.3$ & $0.0074$ & $0.0461$ \\
\bottomrule
\end{tabular}
\par\smallskip
\begin{minipage}{\textwidth}\footnotesize
The VQC/MLP cutoff-$0.3$ rows reuse their E1 seed-$0$ records. Other
cutoff runs use the ablation launcher's $2$ threads. The reference
cutoff-$0.3$ run ends at $1060$ L-BFGS closures; the $3$-thread replay
reproduces E1 exactly at $1065$ closures. Supplementary Section~S1 discusses the comparison. The reference stress-error span across its
five $2$-thread cutoff runs is $0.0506$.
\end{minipage}
\end{table}

\subsection{Training, tuning, and computational accounting}

Table~\ref{tab:training} gives the common nominal schedule. Initial learning rates
are selected independently by minimizing $e_T+e_\sigma$ on fixed random
validation points, never the test grid. Complete candidate histories and
selected scores are released (Supplementary Table~S4). The upper boundary
is selected in $11/21$ E1 and $10/12$ E1b cells; one E1 cell selects the
lower boundary (Supplementary Table~S4). These are conditional comparisons
under the executed tuning and training recipe, not evidence that each
architecture attained its best achievable accuracy. Selection after Adam
also need not identify the best subsequent Adam--L-BFGS endpoint.

The E3 label sweep measures production training-label efficiency
conditional on separately selected hyperparameters. Its data-only rates
were selected once at $N_d=256$, using $256$ exact training samples and
$2000$ exact validation samples shared across the five candidate rates
(Table~\ref{tab:training}; Supplementary Table~S4). This fixed tuning
overhead is outside the production $N_d$ axis. Production runs also
compute exact-label validation diagnostics, but these neither drive
updates nor select checkpoints. Thus $N_d$ is not the total exact-label
requirement for developing or evaluating the tuned learner.

Actual L-BFGS closure calls are objective evaluations, distinct from
optimizer iterations, and vary across runs (Supplementary Table~S5).
The timer covers training, periodic validation, and final test-grid
evaluation, excluding model construction, exact-data preparation, and
tuning. E1 uses the shared CPU with the concurrency and thread settings
in Supplementary Table~S3; other launchers differ. Consequently, the
reported durations are implementation-specific shared-machine timings,
not isolated throughput, total campaign cost, or quantum-hardware cost.
E3 timings are contention-affected observed durations, unsuitable for
isolated-throughput comparisons; E1 is retained as the primary cost comparison.
A reference-MLP
replay at E1's thread setting reproduces its numerical history and endpoint;
the different-thread cutoff run diverges only during L-BFGS
(Table~\ref{tab:xmin}; Supplementary Section~S1). Thread-dependent
reduction order amplified by the line search is a plausible explanation,
not a separately tested causal diagnosis.

\subsection{Metrics and diagnostic protocols}

For $F\in\{T,\sigma\}$ and test grid $\mathcal G$,
\begin{equation}
 e_F=\left[\frac{\sum_{z\in\mathcal G}|F_\vartheta(z)-F^*(z)|^2}
 {\sum_{z\in\mathcal G}|F^*(z)|^2}\right]^{1/2},\qquad
 a_F=\max_{z\in\mathcal G}|F_\vartheta(z)-F^*(z)|.
 \label{eq:metrics}
\end{equation}
E1, E1b, and secondary-speed errors use means and sample standard
deviations with denominator $n-1$. E3 instead reports means and both seed
values or their observed range, because each cell has only two seeds
(Table~\ref{tab:e3}); these ranges are not confidence intervals.
All completed runs, including poor fits, are retained. Best seeds are
descriptive minima. No significance or equivalence test is used.

E2 compares the full physics objective with
\begin{equation}
 \mathcal L_D=\langle(T_\vartheta-T^*)^2\rangle_D
 +\langle(\sigma_\vartheta-\sigma^*)^2\rangle_D,\qquad
 \mathcal L_C=\mathcal L_P+\mathcal L_D.
\end{equation}
The selected derivative, sampling, and variance denominator are specified
in Table~\ref{tab:diagnostics}. Multiplying an objective by a constant
multiplies gradient variance by its square; absolute variance increases
must therefore be interpreted with the stated normalization.

\begin{table}[htbp]
\centering
\caption{Diagnostic settings from the released code and executed records.}
\label{tab:diagnostics}
\begin{tabular}{p{0.25\linewidth}p{0.65\linewidth}}
\toprule
Diagnostic & Setting \\
\midrule
E2 circuit scan & $n_q=2,3,4,5,6,7,8,9,10$; $L=2,4,8$;
$200$ random initializations with seeds $0$--$199$ per cell;
exponential pulse, $\gamma=9$ \\
E2 sampled derivative & First-layer trainable $R_Y$ angle on qubit $0$ \\
E2 points & $256$ fixed uniform interior points, shared between objectives;
$64$ boundary times; point-set generator seed $0$ \\
E2 variance & \texttt{numpy.var} with denominator $200$
(\texttt{ddof=0}) \\
E2 fit & Ordinary least squares of $\log_{10}\Var(\partial\mathcal L/
\partial\theta)$ against $n_q$ \\
E1 gradient summaries & Mean gradient norm over the first and last
$50$ Adam steps \\
E4 checkpoints & Lowest-$e_\sigma$ $2$ E1 seeds per pulse/configuration;
listed in Supplementary Table~S10 \\
E4 evaluation & $41\times40$ grid; binomial samples of each Pauli-$Z$
expectation at $N=10^3,10^4,10^5$ shots per input and observable \\
E4 channel & Pauli-error depolarization after every layer on every qubit;
$p=0,10^{-3},10^{-2}$; density-matrix inference; exact controls at each $p$ \\
\bottomrule
\end{tabular}
\end{table}

E4 evaluates selected checkpoints without retraining. Its dispersion is
across trained checkpoints, not independent repeated shot trials. The
per-qubit channel after each layer is
\begin{equation}
 \mathcal E_{p,q}(\rho)=(1-p)\rho+\frac p3
 (X_q\rho X_q+Y_q\rho Y_q+Z_q\rho Z_q).
 \label{eq:channel}
\end{equation}
Here $p$ is total Pauli-error probability. Independent marginal binomial
draws estimate the readouts using $N$ measurements per input and
observable; commuting readouts could instead share hardware bitstrings.
The evaluator uses each checkpoint's run seed for its sampling generator.
It does not estimate repeated-shot uncertainty or test noisy optimization.

\section{Results}

\subsection{Nominal and output-affecting parameter comparisons}

Table~\ref{tab:e1} is the numerical anchor for E1;
Figure~\ref{fig:e1} displays its field errors. The nominally matched MLP
has lower mean temperature and stress errors in every pulse--configuration
cell. The circuit/MLP mean stress-error ratios range from $1.788$ to
$3.645$ (Table~\ref{tab:e1}). These are ratios of means, not paired
significance tests. The larger circuit's recorded runtime is
$171.7$--$228.3$ times the matched MLP's, compared with
$41.7$--$51.8$ for the smaller circuit (Table~\ref{tab:e1}).

\begin{table}[htbp]
\centering
\caption{E1 nominal-count comparisons. Errors and recorded runtime are mean $\pm$ sample standard deviation across five seeds. $P_{\rm nom}/P_{\rm out}$ lists registered and output-affecting parameter counts, respectively; $P_{\rm out}$ is not a proven independent functional dimension. $R_\sigma$ in each VQC row is its mean stress error divided by its nominally matched MLP mean. Runtime includes training and both evaluations under the stated concurrent execution conditions.}
\label{tab:e1}
\scriptsize
\setlength{\tabcolsep}{3pt}
\begin{tabular}{lllrrrrr}
\toprule
Pulse & Config. & Model & $P_{\rm nom}/P_{\rm out}$ & $e_T$ & $e_\sigma$ & Runtime (s) & $R_\sigma$ \\
\midrule
exp & 4q2L & VQC & $44/30$ & $0.4699\pm0.1428$ & $1.1137\pm0.1711$ & $725.1\pm68.2$ & $3.645$ \\
exp & 4q2L & MLP & $42/42$ & $0.3178\pm0.0738$ & $0.3055\pm0.0608$ & $15.7\pm1.4$ & -- \\
exp & 4q2L & Spectral & $40/40$ & $0.4500\pm0.1532$ & $0.7826\pm0.2417$ & $45.0\pm3.2$ & -- \\
exp & 6q4L & VQC & $124/88$ & $0.1779\pm0.0439$ & $0.7608\pm0.1884$ & $4188.9\pm251.0$ & $3.544$ \\
exp & 6q4L & MLP & $114/114$ & $0.0990\pm0.0584$ & $0.2146\pm0.1375$ & $19.7\pm0.5$ & -- \\
exp & 6q4L & Spectral & $116/116$ & $0.2401\pm0.1689$ & $0.5719\pm0.2035$ & $44.1\pm3.7$ & -- \\
exp & -- & MLP-ref & $12802/12802$ & $0.0125\pm0.0029$ & $0.0539\pm0.0076$ & $214.6\pm9.4$ & -- \\
\midrule
step & 4q2L & VQC & $44/30$ & $0.6688\pm0.1857$ & $1.1934\pm0.3074$ & $624.5\pm15.5$ & $2.506$ \\
step & 4q2L & MLP & $42/42$ & $0.3314\pm0.0683$ & $0.4762\pm0.0743$ & $15.0\pm0.1$ & -- \\
step & 4q2L & Spectral & $40/40$ & $0.7012\pm0.0570$ & $0.9526\pm0.0595$ & $42.5\pm3.4$ & -- \\
step & 6q4L & VQC & $124/88$ & $0.2579\pm0.0336$ & $0.8954\pm0.1136$ & $3644.1\pm21.1$ & $2.513$ \\
step & 6q4L & MLP & $114/114$ & $0.2105\pm0.0904$ & $0.3564\pm0.1015$ & $21.2\pm2.8$ & -- \\
step & 6q4L & Spectral & $116/116$ & $0.3405\pm0.1200$ & $0.3825\pm0.0742$ & $47.3\pm3.4$ & -- \\
step & -- & MLP-ref & $12802/12802$ & $0.0221\pm0.0147$ & $0.0517\pm0.0086$ & $212.6\pm7.8$ & -- \\
\midrule
smooth & 4q2L & VQC & $44/30$ & $0.3260\pm0.0582$ & $0.9826\pm0.0958$ & $758.4\pm115.0$ & $1.788$ \\
smooth & 4q2L & MLP & $42/42$ & $0.2220\pm0.0870$ & $0.5496\pm0.3482$ & $14.6\pm1.0$ & -- \\
smooth & 4q2L & Spectral & $40/40$ & $0.6403\pm0.1968$ & $0.9434\pm0.1900$ & $43.1\pm8.7$ & -- \\
smooth & 6q4L & VQC & $124/88$ & $0.1664\pm0.0342$ & $0.5733\pm0.0696$ & $4473.7\pm466.5$ & $3.364$ \\
smooth & 6q4L & MLP & $114/114$ & $0.0519\pm0.0143$ & $0.1704\pm0.0632$ & $19.6\pm2.6$ & -- \\
smooth & 6q4L & Spectral & $116/116$ & $0.2524\pm0.1288$ & $0.6254\pm0.2388$ & $43.3\pm4.1$ & -- \\
smooth & -- & MLP-ref & $12802/12802$ & $0.0058\pm0.0013$ & $0.0296\pm0.0106$ & $192.8\pm10.6$ & -- \\
\bottomrule
\end{tabular}
\end{table}

\begin{figure}[htbp]
\centering\includegraphics[width=\textwidth]{fig_e1_errors.pdf}
\caption{E1 nominal-count comparison at the settings of
Tables~\ref{tab:problem}--\ref{tab:training}. Bars show mean relative
temperature and stress errors; error bars are sample standard deviations
over five seeds. Solid bars represent the larger configuration and hatched
bars the smaller one; the reference MLP is separate. Circuit
$P_{\rm nom}/P_{\rm out}$ is $124/88$ or $44/30$, with classical counts
listed in Table~\ref{tab:e1}. The plot does not show E1b.}
\label{fig:e1}
\end{figure}

E1b addresses the inert-coordinate imbalance by comparing the same circuits
with retuned smaller classical networks. The MLP again has lower mean
errors for both fields in all six cells, with mean stress-error ratios
$1.558$--$4.493$ (Table~\ref{tab:e1b}). The nominal-count result therefore
does not reverse when this imbalance is reduced. This does not establish
equal intrinsic model dimension or a universal classical advantage.
The spectral baseline is less consistent: its temperature error exceeds
the circuit's in five of six E1b cells, and its larger smooth-pulse stress
error also exceeds the circuit's (Tables~\ref{tab:e1} and \ref{tab:e1b}).

\begin{table}[htbp]
\centering
\caption{E1b classical comparisons targeting circuit output-affecting counts $30$ and $88$, for circuits with $P_{\rm nom}/P_{\rm out}=44/30$ and $124/88$. Each error is mean $\pm$ sample standard deviation over five seeds, with learning rates reselected per cell. $\Delta P$ is the classical count deviation from the target in percent; $R_\sigma$ is the corresponding E1 VQC mean stress error divided by the classical mean.}
\label{tab:e1b}
\scriptsize
\setlength{\tabcolsep}{3pt}
\begin{tabular}{lllrrrrr}
\toprule
Pulse & Target & Model & $P_{\rm nom}/P_{\rm out}$ & $\Delta P$ (\%) & $e_T$ & $e_\sigma$ & $R_\sigma$ \\
\midrule
exp & 4q2L & MLP & $29/29$ & $-3.33$ & $0.3322\pm0.0579$ & $0.4396\pm0.1554$ & $2.533$ \\
exp & 4q2L & Spectral$^{\dagger}$ & $40/40$ & $+33.33$ & $0.4500\pm0.1532$ & $0.7826\pm0.2417$ & $1.423$ \\
exp & 6q4L & MLP & $93/93$ & $+5.68$ & $0.1379\pm0.0302$ & $0.2857\pm0.0607$ & $2.663$ \\
exp & 6q4L & Spectral & $97/97$ & $+10.23$ & $0.3162\pm0.1474$ & $0.6332\pm0.0497$ & $1.201$ \\
\midrule
step & 4q2L & MLP & $29/29$ & $-3.33$ & $0.5283\pm0.1594$ & $0.7658\pm0.3048$ & $1.558$ \\
step & 4q2L & Spectral$^{\dagger}$ & $40/40$ & $+33.33$ & $0.7012\pm0.0570$ & $0.9526\pm0.0595$ & $1.253$ \\
step & 6q4L & MLP & $93/93$ & $+5.68$ & $0.1710\pm0.0267$ & $0.1993\pm0.0391$ & $4.493$ \\
step & 6q4L & Spectral & $97/97$ & $+10.23$ & $0.3425\pm0.1174$ & $0.7213\pm0.2003$ & $1.241$ \\
\midrule
smooth & 4q2L & MLP & $29/29$ & $-3.33$ & $0.2652\pm0.0251$ & $0.4884\pm0.0163$ & $2.012$ \\
smooth & 4q2L & Spectral$^{\dagger}$ & $40/40$ & $+33.33$ & $0.6403\pm0.1968$ & $0.9434\pm0.1900$ & $1.042$ \\
smooth & 6q4L & MLP & $93/93$ & $+5.68$ & $0.0902\pm0.0258$ & $0.3631\pm0.2001$ & $1.579$ \\
smooth & 6q4L & Spectral & $97/97$ & $+10.23$ & $0.3688\pm0.1425$ & $0.8194\pm0.1841$ & $0.700$ \\
\bottomrule
\end{tabular}
\par\smallskip
\begin{minipage}{\textwidth}\footnotesize
$\dagger$ The implemented spectral width search starts at $h=2$, giving a minimum of $40$ parameters. These rows repeat the E1 network outcomes and are not valid output-count matches. A width-one spectral network would have $21$ parameters but was not included in the implemented search. The $97$-parameter spectral model differs from target $88$ by $10.23\%$, slightly exceeding $10\%$.
\end{minipage}
\end{table}

Figure~\ref{fig:fields} shows representative best-stress-seed fields.
The reference MLP attains substantially smaller errors under this recipe
(Table~\ref{tab:e1}); the images illustrate spatial discrepancies without
establishing representational impossibility for any family.

\begin{figure}[htbp]
\centering\includegraphics[width=\textwidth]{fig_fields_exp.pdf}
\caption{Exact and learned exponential-pulse fields at $\gamma=9$.
Temperature is above stress; columns show the exact solution, the larger
circuit ($P_{\rm nom}/P_{\rm out}=124/88$), nominally matched MLP
($114/114$), and reference MLP ($12802/12802$). Learned examples use the
minimum-stress-error E1 seeds, identified in Supplementary Table~S6.
Quantitative color bars give dimensionless $T$ and $\sigma$. The
characteristic $x=3t$ follows the wave speed in Table~\ref{tab:problem}.
These are selected examples, not replicated error summaries.}
\label{fig:fields}
\end{figure}

\subsection{Gradient variance}

Figure~\ref{fig:e2} and Table~\ref{tab:e2summary} summarize the variance
of a selected derivative of the full physics objective. The physics/data
variance multipliers range from $5.19$--$20.31$, $70.89$--$375.41$, and
$2726.20$--$9311.63$ across the respective depths (Supplementary Table~S9).
The deepest whole-range slopes are negative, but its largest-qubit refits
are positive (Table~\ref{tab:e2summary}). Thus neither unchanged decay
nor persistent exponential scaling is established.

\begin{table}[htbp]
\centering
\caption{Descriptive ordinary-least-squares slopes of $\log_{10}\operatorname{Var}(\partial\mathcal L/\partial\theta)$ versus qubit count. The selected derivative is the first-layer $R_Y$ angle on qubit $0$; each variance uses $200$ initializations. The final row refits the largest-qubit portion of the deepest scan. No slope-equivalence test or asymptotic exponent is inferred.}
\label{tab:e2summary}
\footnotesize
\setlength{\tabcolsep}{3pt}
\begin{tabular}{rrrrr}
\toprule
Depth $L$ & Fitted $n_q$ range & Data & Full physics & Combined \\
\midrule
2 & $2$--$10$ & $-0.008$ & $+0.027$ & $+0.020$ \\
4 & $2$--$10$ & $-0.038$ & $-0.076$ & $-0.074$ \\
8 & $2$--$10$ & $-0.206$ & $-0.227$ & $-0.227$ \\
8 & $7$--$10$ & $+0.014$ & $+0.020$ & $+0.020$ \\
\bottomrule
\end{tabular}
\end{table}

\begin{figure}[htbp]
\centering\includegraphics[width=\textwidth]{fig_e2_gradvar.pdf}
\caption{Variance of the first-layer measured-qubit $R_Y$ derivative at
random initialization, for the objectives and sampling in
Table~\ref{tab:diagnostics}. Integer qubit ticks and distinct markers
separate the full physics and combined objectives. The curves use
population-denominator variance over $200$ initializations per cell;
they are not mean--standard-deviation bands. Full values are in
Supplementary Table~S8 and descriptive slopes in
Table~\ref{tab:e2summary}.}
\label{fig:e2}
\end{figure}

At fixed depth a local readout has a bounded backward causal cone, so
indefinite qubit growth cannot be assumed to reproduce a whole-range
finite-size fit \citep{Cerezo2021}. This observation limits extrapolation;
it does not demonstrate that cone saturation caused the observed tail,
whose circuit can still span the tested ring. Individual gradient draws
are not retained, and no slope-equivalence test is available. E2 therefore
answers a narrower question than a barren-plateau theorem and does not
isolate the contribution of the PDE residual from the other physics terms.

\subsection{Efficiency}
\label{sec:efficiency}

The completed E3 sweeps address parameter count, collocation budget, and
production training-label budget conditional on separately selected
hyperparameters for the exponential pulse. Table~\ref{tab:e3}
reports means and observed ranges over the declared two seeds;
Supplementary Section~S9 retains every individual run, matching deviation,
learning rate, point budget, runtime, and realized closure count.
The parameter and collocation sweeps use the reference physics recipe;
data-only training uses $\mathcal L_D$, fixed exact interior samples,
separately selected rates, and no boundary, rate, or PDE penalty or L-BFGS
(Table~\ref{tab:e3}; Supplementary Tables~S3--S4).

\paragraph{Parameter count.}
The nominally matched MLP has lower mean stress error at every tested
configuration, by factors $2.0$--$8.5$ (Table~\ref{tab:e3};
Figure~\ref{fig:e3pareto}). Even the best circuit run, $e_\sigma=0.3951$
at $(n_q,L)=(6,8)$, exceeds the smallest MLP's mean $0.3067$ at
$42$ parameters. The best circuit mean is $0.4464$ at
$244/208$ nominal/output-affecting parameters (Table~\ref{tab:e3}).
Changing the horizontal coordinate to the output-affecting count does
not change this accuracy ordering. It does not, however, establish strict
Pareto domination of every circuit point: the shallow circuits have fewer
output-affecting coordinates than the smallest E3 MLP
(Table~\ref{tab:e3}; Supplementary Table~S2). E1b supplies the separate
comparison with smaller classical networks (Table~\ref{tab:e1b}).

The depth-$2$ circuits with $4$, $6$, and $8$ qubits realize the same
$30$-coordinate output family, not merely the same output-affecting count
(Supplementary Section~S2). The final CZ layer cancels in the measured
$Z$ expectations; propagation through the preceding ring leaves only
qubits $0,1$ and their neighbors $2,n_q-1$. Relabeling these active qubits
therefore preserves the family. The supplied execution of
\path{family_equivalence.py} confirms the parameter transplant to
numerical precision (Supplementary Sections~S2 and~S8).
Their stress means $1.0127$, $1.4455$, and $2.0390$ are differently
initialized and trained realizations of that same family with different
simulation cost, not a capacity trend (Table~\ref{tab:e3}). At fixed
width, increasing depth reduces stress error from $1.4455$ through
$0.7352$ to $0.4464$ (Table~\ref{tab:e3}); this remains an observed
trend without an identified mechanism. The deepest circuit
approaches the MLP's temperature accuracy, $0.0598$ versus $0.0557$,
without an equivalence test. Spectral stress errors fall between circuit
and MLP in all but the $(6,4)$ design, where the spectral mean is slightly
larger than the circuit's (Table~\ref{tab:e3}).

The classical rows for circuit targets $(4,4)$ and $(8,2)$ share
architectures and identical numerical outcomes; their repeated executions
are marked in Table~\ref{tab:e3}. Their MLP count deviation is $+10.71\%$;
the $(6,2)$ MLP deviation is $-10.94\%$, both outside the prescribed
tolerance (Supplementary Table~S2). The $(4,2)$ and $(6,4)$ parameter
runs exactly reproduce the corresponding E1 numerical histories and
errors for the same seeds, although recorded runtimes differ
(Supplementary Section~S1). They do not add independent replication.

\begin{figure}[htbp]
\centering\includegraphics[width=0.80\textwidth]{fig_e3_pareto.pdf}
\caption{E3 stress error versus parameter count for the exponential pulse,
$\gamma=9$, using means of seeds $0,1$. Red filled circles use circuit
nominal counts; hollow squares show the same errors at output-affecting
counts. Blue and green points denote nominally matched MLP and spectral
models. Lines connect evaluated architectures and are not fitted Pareto
envelopes. The duplicated classical outcomes at the common nominal target
are shared. The depth-$2$ circuits at $4$, $6$, and $8$ qubits are
different training realizations of one $30$-coordinate output family,
as established by the causal structure and supplied parameter-transplant
check (Supplementary Sections~S2 and~S8); their width sequence is not a
capacity comparison. Table~\ref{tab:e3} gives both-seed ranges and Supplementary
Table~S2 gives every count and matching deviation.}
\label{fig:e3pareto}
\end{figure}

\paragraph{Collocation budget.}
Increasing the number of residual points does not systematically lower
either family's errors or close the stress gap in the tested range
(Table~\ref{tab:e3}; Figure~\ref{fig:e3colloc}). Across shared budgets,
the circuit/MLP stress-mean ratio is $2.0$--$4.4$, and every completed
circuit run has stress error above $0.5$ (Table~\ref{tab:e3}).
The nonmonotone means and limited seed replication do not prove that
either architecture is generally free of collocation limitations.

The circuit uses $N_c=64,256,512,1024$; the MLP also uses $4096$.
The baseline at $512$ reuses the parameter-sweep runs. The circuit jobs
at $4096$ were dropped after an initial execution was stopped for cost;
the measured-rate projection was about $32400\,\mathrm{s}$ per seed
after $2100\,\mathrm{s}$ of observation, including a fixed L-BFGS set
of $16384$ points (Table~\ref{tab:e3}; Supplementary Table~S3).
This is a declared deviation from the planned sweep, not an unreported
completed outcome. The retained campaign contains $66$ of the planned
$68$ runs (Supplementary Table~S1).

\begin{figure}[htbp]
\centering\includegraphics[width=0.80\textwidth]{fig_e3_colloc.pdf}
\caption{E3 stress error versus collocation points per Adam step.
Physics-only training uses fixed L-BFGS sets of $4N_c$ points.
Means use seeds $0,1$ for the $(6,4)$ circuit ($124/88$ parameters)
and nominally matched MLP ($114/114$). Circuit budgets are
$64,256,512,1024$; the MLP also includes $4096$. The circuit's largest
planned budget was dropped for cost. The baseline at $512$ is reused
from the parameter sweep. Observed ranges are in Table~\ref{tab:e3}.}
\label{fig:e3colloc}
\end{figure}

\paragraph{Production training labels.}
The $N_d$ axis counts exact interior samples used for production updates,
conditional on the separate tuning at $N_d=256$. It excludes the fixed
tuning set of $256$ exact training samples plus $2000$ exact validation
samples shared across the five candidate rates, and the production
validation diagnostics (Table~\ref{tab:training}; Supplementary Table~S4).
Those diagnostics neither drive updates nor select checkpoints.
Under supervision, circuit stress means fall from $2.1382$ at $16$ labels
to $0.2918$ at $256$ and $0.2178$ at $1024$; the last value is
$3.375$ times lower than its physics-only mean for the same seeds
(Table~\ref{tab:e3}; Figure~\ref{fig:e3data}). Thus the same architecture
can represent a more accurate stress field than physics training attained.
The gap is partly training-objective dependent, rather than evidence of
representational impossibility. This comparison changes information,
objective, rate, and optimizer schedule, so it cannot establish that the
loss landscape is the main cause.

The MLP remains more accurate at every label budget, with stress-mean
ratios $1.4$--$3.3$. Its means change little from $0.0884$ to $0.0872$
between the largest budgets. The circuit at $1024$ labels still exceeds
the MLP mean at $64$ labels, $0.2178$ versus $0.1903$
(Table~\ref{tab:e3}). These observed production training budgets differ by
a factor of $16$; they estimate neither an equal-error sample-complexity
ratio nor a total label-acquisition ratio. At the smallest
budget both families have mean errors above unity for both fields
(Table~\ref{tab:e3}).

\begin{figure}[htbp]
\centering\includegraphics[width=0.80\textwidth]{fig_e3_data.pdf}
\caption{E3 stress error versus fixed exact interior training samples
$N_d$ in production runs conditional on separately selected hyperparameters.
Means use seeds $0,1$ for the same circuit and MLP
as Figure~\ref{fig:e3colloc}. Training uses $2000$ Adam steps, no PDE,
boundary, or initial-rate penalty, and no L-BFGS. Validation-selected
rates are $0.03$ for the circuit and $0.1$ for the MLP, selected once at
$N_d=256$ using $256$ training and $2000$ validation samples shared by
the five candidates. The horizontal axis excludes these tuning labels
and production validation diagnostics, which neither drive updates nor
select checkpoints (Table~\ref{tab:training}). Table~\ref{tab:e3}
provides the two-seed ranges.}
\label{fig:e3data}
\end{figure}

E3 timings are contention-affected observed durations, unsuitable for
isolated-throughput comparisons.
For example, the deepest physics circuit takes $8683$--$11552\,\mathrm{s}$,
whereas its MLP takes $19$--$22\,\mathrm{s}$; the supervised circuit ranges
from $28$ to $156\,\mathrm{s}$ across label budgets (Supplementary
Table~S16). These are not controlled cost ratios. E1 remains the primary
wall-clock comparison (Table~\ref{tab:e1}; Supplementary Table~S3).

\begin{table}[!htbp]
\centering\footnotesize
\setlength{\tabcolsep}{4pt}
\caption{E3 errors: mean [minimum, maximum] over seeds $0,1$ ($n=2$).
$P$ denotes nominal/output-affecting counts. Parameter rows use $N_c=512$;
collocation rows vary $N_c$ with $4N_c$ fixed L-BFGS points; data rows vary
$N_d$ conditional on separate tuning (Table~\ref{tab:training}) and omit
all physics losses and L-BFGS. All physics rates are $0.1$;
data rates are $0.03$ (VQC) and $0.1$ (MLP). Budgets and per-seed times,
closures, and matching deviations are retained in the supplement.
The $N_c=512$ rows reuse the parameter-sweep records.
A dagger marks classical outcomes identical to the $(4,4)$ executions,
not additional independent evidence.}\label{tab:e3}
\begin{tabular}{llcll}
\toprule
Design/budget & Model & $P$ & $e_T$ & $e_\sigma$ \\
\midrule
\multicolumn{5}{l}{\textit{Parameter sweep: $(n_q,L)$}}\\
$4,2$ & VQC & $44/30$ & $0.4030\,[0.2561,0.5500]$ & $1.0127\,[0.8460,1.1794]$ \\
$4,2$ & MLP & $42/42$ & $0.2755\,[0.2723,0.2788]$ & $0.3067\,[0.2311,0.3823]$ \\
$4,2$ & Spectral & $40/40$ & $0.4056\,[0.4001,0.4112]$ & $0.8616\,[0.7746,0.9487]$ \\
$4,4$ & VQC & $84/70$ & $0.1947\,[0.1905,0.1990]$ & $0.9752\,[0.9030,1.0473]$ \\
$4,4$ & MLP & $93/93$ & $0.1052\,[0.1005,0.1100]$ & $0.2403\,[0.2196,0.2610]$ \\
$4,4$ & Spectral & $78/78$ & $0.2830\,[0.1796,0.3864]$ & $0.5552\,[0.4204,0.6899]$ \\
$6,2$ & VQC & $64/30$ & $0.6330\,[0.6162,0.6498]$ & $1.4455\,[1.3497,1.5413]$ \\
$6,2$ & MLP & $57/57$ & $0.1530\,[0.1020,0.2040]$ & $0.3598\,[0.2854,0.4341]$ \\
$6,2$ & Spectral & $59/59$ & $0.3909\,[0.2988,0.4830]$ & $0.8662\,[0.8510,0.8815]$ \\
$6,4$ & VQC & $124/88$ & $0.1657\,[0.1376,0.1939]$ & $0.7352\,[0.5919,0.8785]$ \\
$6,4$ & MLP & $114/114$ & $0.1402\,[0.0844,0.1959]$ & $0.3640\,[0.3409,0.3872]$ \\
$6,4$ & Spectral & $116/116$ & $0.2159\,[0.1650,0.2667]$ & $0.7825\,[0.7594,0.8056]$ \\
$6,8$ & VQC & $244/208$ & $0.0598\,[0.0591,0.0606]$ & $0.4464\,[0.3951,0.4977]$ \\
$6,8$ & MLP & $249/249$ & $0.0557\,[0.0359,0.0756]$ & $0.1195\,[0.0935,0.1455]$ \\
$6,8$ & Spectral & $249/249$ & $0.1420\,[0.1327,0.1513]$ & $0.3823\,[0.3583,0.4064]$ \\
$8,2$ & VQC & $84/30$ & $0.8757\,[0.6694,1.0819]$ & $2.0390\,[1.7132,2.3649]$ \\
$8,2$ & MLP$^{\dagger}$ & $93/93$ & $0.1052\,[0.1005,0.1100]$ & $0.2403\,[0.2196,0.2610]$ \\
$8,2$ & Spectral$^{\dagger}$ & $78/78$ & $0.2830\,[0.1796,0.3864]$ & $0.5552\,[0.4204,0.6899]$ \\
\midrule\multicolumn{5}{l}{\textit{Collocation sweep: $N_c$, $(n_q,L)=(6,4)$}}\\
$64$ & VQC & $124/88$ & $0.2199\,[0.1051,0.3348]$ & $0.8198\,[0.8078,0.8318]$ \\
$64$ & MLP & $114/114$ & $0.1287\,[0.0870,0.1705]$ & $0.2041\,[0.2023,0.2060]$ \\
$256$ & VQC & $124/88$ & $0.2158\,[0.2140,0.2176]$ & $0.6687\,[0.5376,0.7998]$ \\
$256$ & MLP & $114/114$ & $0.1198\,[0.0906,0.1491]$ & $0.2575\,[0.2438,0.2713]$ \\
$512$ & VQC & $124/88$ & $0.1657\,[0.1376,0.1939]$ & $0.7352\,[0.5919,0.8785]$ \\
$512$ & MLP & $114/114$ & $0.1402\,[0.0844,0.1959]$ & $0.3640\,[0.3409,0.3872]$ \\
$1024$ & VQC & $124/88$ & $0.2139\,[0.1815,0.2463]$ & $0.7594\,[0.6938,0.8250]$ \\
$1024$ & MLP & $114/114$ & $0.0962\,[0.0770,0.1154]$ & $0.1726\,[0.0886,0.2566]$ \\
$4096$ & MLP & $114/114$ & $0.1165\,[0.1118,0.1213]$ & $0.2968\,[0.2496,0.3441]$ \\
\midrule\multicolumn{5}{l}{\textit{Label sweep: $N_d$, $(n_q,L)=(6,4)$}}\\
$16$ & VQC & $124/88$ & $1.6397\,[0.9935,2.2859]$ & $2.1382\,[1.7633,2.5130]$ \\
$16$ & MLP & $114/114$ & $1.0870\,[0.7236,1.4504]$ & $1.5074\,[1.2370,1.7778]$ \\
$64$ & VQC & $124/88$ & $0.2798\,[0.2571,0.3024]$ & $0.5336\,[0.4412,0.6260]$ \\
$64$ & MLP & $114/114$ & $0.0802\,[0.0738,0.0865]$ & $0.1903\,[0.1577,0.2229]$ \\
$256$ & VQC & $124/88$ & $0.1907\,[0.1873,0.1940]$ & $0.2918\,[0.2882,0.2953]$ \\
$256$ & MLP & $114/114$ & $0.1007\,[0.0873,0.1140]$ & $0.0884\,[0.0773,0.0995]$ \\
$1024$ & VQC & $124/88$ & $0.1920\,[0.1769,0.2070]$ & $0.2178\,[0.2146,0.2211]$ \\
$1024$ & MLP & $114/114$ & $0.0776\,[0.0682,0.0870]$ & $0.0872\,[0.0858,0.0885]$ \\
\bottomrule
\end{tabular}
\end{table}
\FloatBarrier

\subsection{Noise and readout sensitivity}

Table~\ref{tab:e4ratios} compares each selected checkpoint pair with its
same-grid, zero-noise exact-expectation control. Supplementary Table~S10
merges seed identities, individual gains, offsets, and clean errors;
Tables~S11--S12 preserve every noise condition. These controls use the
noise grid, not the E1-grid metadata stored alongside the records.

For the larger exponential-pulse circuit, mean stress error is
$0.5804\pm0.0166$ clean and $2.5758\pm0.2728$, $0.9735\pm0.0950$, and
$0.6299\pm0.0251$ at $10^3$, $10^4$, and $10^5$ shots per input and
observable, respectively ($n=2$; Supplementary Table~S11). At the lowest
budget the larger-circuit stress multipliers span $1.642$--$4.438$,
whereas smaller-circuit multipliers span $8.637$--$22.396$
(Table~\ref{tab:e4ratios}).

\begin{table}[htbp]
\centering
\caption{Ratios of corrected E4 mean stress error to the corresponding same-grid clean mean. Shot columns use $p=0$; channel columns use exact expectations. Ratios are descriptive ratios of two-seed means, not means of paired ratios or significance tests.}
\label{tab:e4ratios}
\small
\setlength{\tabcolsep}{4pt}
\begin{tabular}{llrrrrr}
\toprule
Pulse & $P_{\rm nom}/P_{\rm out}$ & $N=10^3$ & $10^4$ & $10^5$ & $p=10^{-3}$ & $10^{-2}$ \\
\midrule
exp & $44/30$ & $22.396$ & $7.349$ & $2.600$ & $0.998$ & $0.982$ \\
exp & $124/88$ & $4.438$ & $1.677$ & $1.085$ & $0.991$ & $1.057$ \\
step & $44/30$ & $8.637$ & $2.929$ & $1.375$ & $0.995$ & $0.978$ \\
step & $124/88$ & $1.642$ & $1.079$ & $1.009$ & $0.988$ & $0.936$ \\
smooth & $44/30$ & $9.540$ & $3.307$ & $1.510$ & $1.000$ & $1.003$ \\
smooth & $124/88$ & $2.730$ & $1.267$ & $1.027$ & $1.003$ & $1.058$ \\
\bottomrule
\end{tabular}
\end{table}

Exact-expectation depolarization yields larger exponential-pulse stress
errors of $0.5753\pm0.0184$ at $p=10^{-3}$ and $0.6135\pm0.1210$ at
$p=10^{-2}$ ($n=2$; Supplementary Table~S11). Across all cells the
stress changes are modest (Table~\ref{tab:e4ratios}), not several-fold
channel penalties. Small decreases do not establish beneficial noise.
Temperature can behave differently: the smaller smooth-pulse circuit
changes from $0.3249\pm0.0189$ clean to $1.2674\pm1.3552$ at
$p=10^{-2}$ ($n=2$; Supplementary Table~S11).

For a readout $z=\langle Z\rangle$ and $F=t(\alpha z+\beta)$,
independent measurement statistics give
\begin{equation}
 \Var(\widehat F\mid z)=\frac{t^2\alpha^2(1-z^2)}N,
 \qquad \delta F=t\alpha\,\delta z.
 \label{eq:amplification}
\end{equation}
Individual trained gains vary strongly even within the selected pairs
(Supplementary Table~S10). Large gains amplify fluctuations; offsets do
not cancel them. Channel-induced bias also depends on the actual change
in $z$ and the existing approximation error, explaining why large gains
alone do not imply large depolarizing penalties. Selection by best clean
E1 stress error limits these findings to the evaluated checkpoints.

\subsection{Secondary wave speeds}
\label{sec:secondary}

Table~\ref{tab:gamma} reports the completed higher-wave-speed experiments
and the same-seed subset of E1. The MLP retains lower mean temperature
and stress errors than the circuit. Circuit stress error increases from
$0.8274\pm0.2146$ at the original wave-speed parameter to
$1.1971\pm0.4658$ and $1.4485\pm0.8240$ at the secondary values
($n=3$; Table~\ref{tab:gamma}). Its largest individual error is $2.3413$
(Supplementary Table~S17). The growing dispersion cautions against
interpreting only the means.

Every family has larger mean errors at the secondary speeds than in the
full five-seed E1 table (Tables~\ref{tab:e1} and \ref{tab:gamma}), but this
statement changes the seed population.
With the same seeds, the MLP's stress mean improves from $0.2779$ to
$0.2353$ at the intermediate speed (Table~\ref{tab:gamma}).
Consequently, the evidence supports a preserved circuit/MLP ordering,
not uniform degradation for every family and seed subset.

The wave front reaches the outer boundary at $t=6/\sqrt\gamma$:
$2$, $1.2$, and approximately $0.8571$ for the three tested values
(Tables~\ref{tab:problem} and \ref{tab:gamma}). Exact temperature is
independent of $\gamma$, so its changed learning errors reflect the coupled
training problem. Increasing $\gamma$ changes the stress normalization,
initial amplitude, and relative residual scaling. In particular,
Eq.~\eqref{eq:loss} weights the squared unscaled stress-rate mismatch by
$\gamma^{-2}$; the normalized rate target remains unchanged.
Earlier front exit and changed loss scaling are relevant differences,
not separately isolated causes of the observed errors. Rates were
transferred without retuning, and timing shares the resource caveats
of the overlapping campaign (Supplementary Tables~S3 and~S17).

\begin{table}[htbp]
\centering\small
\caption{Secondary wave speeds: mean $\pm$ sample standard deviation of
errors over seeds $0,1,2$ ($n=3$). The $\gamma=9$ rows are the same-seed
subset of E1, not new runs. All runs use the exponential pulse and the
E1 physics recipe; rates transfer from $\gamma=9$: $0.1$, except $0.01$
for the reference. Circuit, MLP, spectral, and reference counts are
$124/88$, $114/114$, $116/116$, and $12802/12802$, respectively.
The reference is not parameter matched. Per-seed errors, times, and
closure counts are retained in the supplement.}
\label{tab:gamma}
\begin{tabular}{rlll}
\toprule
$\gamma$ & Model & $e_T$ & $e_\sigma$ \\
\midrule
$9$ & VQC & $0.1910\pm0.0520$ & $0.8274\pm0.2146$ \\
$9$ & MLP & $0.1164\pm0.0693$ & $0.2779\pm0.1511$ \\
$9$ & Spectral & $0.1895\pm0.0683$ & $0.6309\pm0.2636$ \\
$9$ & Reference & $0.0117\pm0.0038$ & $0.0488\pm0.0044$ \\
\midrule
$25$ & VQC & $0.2927\pm0.0458$ & $1.1971\pm0.4658$ \\
$25$ & MLP & $0.1279\pm0.0525$ & $0.2353\pm0.1525$ \\
$25$ & Spectral & $0.3180\pm0.0970$ & $0.6342\pm0.3302$ \\
$25$ & Reference & $0.0221\pm0.0028$ & $0.1044\pm0.0357$ \\
\midrule
$49$ & VQC & $0.3868\pm0.1884$ & $1.4485\pm0.8240$ \\
$49$ & MLP & $0.2468\pm0.0621$ & $0.4511\pm0.1311$ \\
$49$ & Spectral & $0.4400\pm0.1898$ & $0.6894\pm0.2043$ \\
$49$ & Reference & $0.0421\pm0.0014$ & $0.0956\pm0.0208$ \\
\bottomrule
\end{tabular}
\end{table}

\section{Discussion}

\subsection{What the comparisons establish}

RQ1 has a negative descriptive answer under the executed recipe: the MLP
has lower mean field errors in the nominal and output-affecting comparisons
(Tables~\ref{tab:e1} and \ref{tab:e1b}). Output-affecting matching addresses
known inert coordinates, while leaving capacity and optimization
differences. The result neither invalidates favorable quantum-PINN
comparisons on other equations and embeddings \citep{Berger2025} nor
establishes an optimal classical or quantum architecture. Search-boundary
choices and single-seed tuning further limit generalization.

RQ2 supports increased scale of one full-objective gradient statistic,
without an asymptotic trainability conclusion
(Table~\ref{tab:e2summary}; Supplementary Table~S9). Its fixed readouts and input derivatives
differ from qubit-indexed PDE discretizations \citep{Ukwatta2026}.
RQ3 has a negative answer relative to the MLP over the tested budgets:
no stress-efficiency advantage appears in parameter, collocation, or
production training-label sweeps conditional on separately selected
hyperparameters (Table~\ref{tab:e3}). The label axis excludes the fixed
tuning overhead and production validation diagnostics; it does not measure
total label-acquisition efficiency (Table~\ref{tab:training}). The shallow
width sequence repeats one output family with different initialization,
training, and simulation cost (Supplementary Section~S2).
The spectral model is a weaker stress control
and does not beat the circuit at every parameter target; the strongest
negative comparison is with the MLP. Near equality of the deepest
circuit's temperature mean does not establish parity for the coupled
task. Supervised training demonstrates better attainable circuit stress
accuracy, but the classical gap persists. These findings are restricted
to the sampled architectures, transferred physics rates, and two seeds
(Table~\ref{tab:e3}; Supplementary Table~S2). E1 provides the less
contention-affected cost comparison; E3 timings cannot establish isolated
compute scaling or project hardware speedups.

\subsection{Representation, weighting, and classical alternatives}

The spectral comparator tests a particular random basis, not a classical
surrogate constructed from the trained circuit. It therefore cannot
establish or refute dequantization \citep{Schreiber2023,Sahebi2025}.
Trainable circuit encoding scales and fixed sampled frequencies differ
substantially. More targeted spectral controls, learned frequencies, or
surrogates of the fitted circuit would address different questions.

Wave fronts and multiple temporal scales make frequency content relevant;
Fourier-feature analyses of classical PINNs provide context
\citep{WangFourier2021}. Unequal convergence of loss components is another
documented PINN difficulty \citep{WangNTK2022}. Neither mechanism is
diagnosed by the present field plots or aggregate losses. Adaptive
weighting \citep{Setty2025} and alternative quantum embeddings
\citep{TranEmbedding2026} merit controlled tests, rather than attribution
of the observed errors to a single presumed representational limit.
The supervised E3 result strengthens this distinction: the circuit's
physics-only error is not a lower bound on its representation accuracy
(Table~\ref{tab:e3}). It does not identify how much of the deficit arises
from loss weighting, optimization, information access, or the initial-rate
relaxation.

\subsection{Readout amplification and limitations}

Equation~\eqref{eq:amplification} makes unconstrained gains a practical
inference concern even after exact-expectation training. Gain penalties,
variance regularization, or shot-aware training could trade clean accuracy
against sampling cost \citep{Kreplin2024}. E4 contains no intervention
testing those remedies. Its fixed-checkpoint inference experiment is also
distinct from noisy-gradient barren-plateau results \citep{WangNoise2021}.

The study uses compact circuits and classical networks, far below many
practical PINN sizes, simulated statevectors, and no hardware runs.
Common iteration caps do not equalize derivative work or realized closure
calls. Exact outer-boundary and validation information is available to
every physics-trained family. The initial-rate strip leaves the continuum
stress problem underdetermined. Its nonmonotone single-seed sensitivity
differs across models: the declared cutoff gives the best tested stress
outcome for the circuit and reference MLP but the worst for the matched
MLP (Table~\ref{tab:xmin}). It is therefore neither a generally validated
accuracy improvement nor a substitute for complete initial data.
The dropped high-collocation circuit jobs, sparse efficiency replication,
thread-sensitive reference replay, and overlapping execution conditions
further limit scalability and timing conclusions (Supplementary
Sections~S1 and~S9).
The free-surface startup corner and rectangular switch-off interface also
limit smooth strong-form approximation. Verification is pointwise at a
single finite-difference resolution, not a uniform accuracy bound.
The gradient study samples one derivative, while E4 selects clean-error
minima and lacks repeated shot realizations. These restrictions define the
scope of the negative comparison.

\section{Conclusions}

Under the declared tuning, sampling, and initial-rate recipe, the classical
MLP achieves lower mean temperature and stress errors than the tested
circuit in both nominal and output-affecting count comparisons
(Tables~\ref{tab:e1} and \ref{tab:e1b}). The full physics objective raises
a selected gradient variance, but its finite-size slopes do not establish
barren-plateau mitigation (Table~\ref{tab:e2summary}; Supplementary Table~S9).
The completed parameter, collocation, and production training-label
sweeps show no circuit stress advantage over the MLP. The label comparison
is conditional on separately selected hyperparameters and excludes tuning
and diagnostic labels (Table~\ref{tab:training}). Supervision reveals better attainable circuit accuracy
than the physics recipe achieved (Table~\ref{tab:e3}). Secondary speeds
preserve the circuit/MLP ordering, with larger circuit stress dispersion
(Table~\ref{tab:gamma}). Trained readout gains
produce substantial finite-shot sensitivity in selected circuits, whereas
the tested exact-expectation depolarization changes their stress errors
modestly (Table~\ref{tab:e4ratios}).

For laser processing, thermal-stress control, and thermal barrier coating
assessment, the practical lesson is to verify stress initial data,
normalization, field accuracy, and computational cost together. The
corrected reference fields and explicit conversion to pascals support
that verification. Extension to realistic materials requires complete
initial-rate treatment and richer thermal-mechanical physics; this
benchmark establishes neither quantum advantage nor a calibrated
engineering design tool.

%TC:ignore
\section*{Data and code availability}
The assembled archive \path{QPINN_LTE_release_v1_2026-10-21.zip}
contains the analysis and reference code, all launchers and generators,
parameter-audit and family-equivalence scripts, and the four verification
scripts. Its \path{results/} directory contains $464$ per-run JSON records:
the $440$ retained in \path{runs_all.json} and segregated superseded
launches used by no table. It also includes $456$ checkpoint state-dictionary
files (\texttt{.pt}, $7.8$ MB), E2 aggregate records
(\path{e2_exp_g9.json}), and corrected checkpoint-labeled E4 records
(\path{e4_noise.json}). Execution logs in \path{logs/} ($2.2$ MB)
include queue dispatchers and the initial-rate weight-ablation logs and
variants. The package contains results package v6, tables and figures,
histories and closure counts, verification and audit outputs, the
family-equivalence check, checkpoint gains, and selected rates
(Supplementary Section~S1 and Table~S1). The \path{RELEASE_README}
maps every manuscript table and figure to its inputs and generator and
records regeneration commands and launcher thread settings; Supplementary
Sections~S8--S9 give reconstruction details and individual E3 and
secondary-speed outcomes. Repository address or DOI: to be inserted by
the author at submission.
% Author to insert and confirm the accessible repository address or DOI at submission.

\section*{CRediT}
Hussain Al-Qahtani: Conceptualization, Methodology, Software, Validation,
Formal analysis, Investigation, Writing---original draft,
Writing---review and editing.

\section*{Declaration of competing interests}
The author declares no competing interests.
% to be confirmed by the author before submission

\section*{Acknowledgements}
The author acknowledges King Fahd University of Petroleum and Minerals
(KFUPM) for institutional support. Generative AI assisted manuscript
organization, language editing, and consistency checks. The author is
responsible for the scientific content, source verification, and final text.

\appendix
\section{Corrected closed-form solution and verification}
\label{app:exact}

The expressions below define the exact reference; Supplementary
Section~S7 gives the correction inventory and executed pointwise checks.
The transformed initial-rate terms must be retained before cancellation.

\subsection{Transform solution and the initial stress rate}

Let $\widehat f(s)=\mathcal L[f](s)$, $a=\sqrt\gamma$, and
$\tau=x/a$, using the branch of $\sqrt s$ positive on the positive real axis.
Zero initial fields and
$T_t(x,0^+)=e^{-x}f(0^+)$ together with
$\sigma_t(x,0^+)=-\gamma e^{-x}f(0^+)$ give
\begin{equation}
 \widehat\sigma_{xx}-\frac{s^2}{\gamma}\widehat\sigma
 =s^2\widehat T.
\end{equation}
The initial-rate terms on the two sides cancel. Solving the spatial boundary
value problems gives
\begin{align}
 \widehat T(x,s)
 &=\widehat f(s)\left[
 \frac{e^{-x}}{s-1}-\frac{e^{-x\sqrt s}}{\sqrt s(s-1)}\right],\\
 \widehat\sigma(x,s)
 &=\widehat f(s)\left[
 \mathcal{A}(s)\bigl(e^{-x}-e^{-s\tau}\bigr)
 -\mathcal{B}(s)\bigl(e^{-x\sqrt s}-e^{-s\tau}\bigr)\right],\\
 \mathcal{A}(s)&=\frac{\gamma s^2}{(s-1)(\gamma-s^2)},\qquad
 \mathcal{B}(s)=\frac{\gamma\sqrt s}{(s-1)(\gamma-s)}.
\end{align}
The apparent positive poles at $s=1$, $s=a$, and $s=\gamma$ cancel in
the complete expression, so they do not produce growing physical modes.
At $s=1$ the cancellation is between terms; at $s=a$ and $s=\gamma$
the corresponding exponential differences vanish. The analogous assertion
that the first difference vanishes at $s=-a$ would be incorrect.

\subsection{Exponential response and pulse superposition}

Let $\mathcal T_\beta$ and $\mathcal S_\beta$ denote the responses to
$e^{-\beta t}H(t)$, with unit amplitude. For $t>0$, define
\begin{align}
 \xi&=\frac{x}{2\sqrt t},\qquad \eta=\sqrt{\beta t},\\
 J_\beta(x,t)&=e^{-\beta t}
  \operatorname{Im}\left[e^{i\sqrt\beta x}
     \operatorname{erfc}(\xi+i\eta)\right],\\
 C_1&=\frac{\gamma}{(\gamma-1)(1+\beta)},&
 C_2&=\frac{\gamma}{(\gamma-1)(\gamma+\beta)},&
 C_3&=-\frac{\gamma}{(1+\beta)(\gamma+\beta)},\\
 r_+&=-\frac{\gamma a}{2(a-1)(a+\beta)},&
 r_-&=\frac{\gamma a}{2(a+1)(a-\beta)},&
 r_\beta&=\frac{\gamma\beta^2}{(1+\beta)(\beta^2-\gamma)}.
\end{align}
For $\beta>0$,
\begin{equation}
 \mathcal T_\beta(x,t)=\frac{1}{1+\beta}\left[
 \frac{e^{t-x}}{2}\operatorname{erfc}(\sqrt t-\xi)
 +\frac{e^{t+x}}{2}\operatorname{erfc}(\sqrt t+\xi)
 -e^{-x-\beta t}-\frac{J_\beta(x,t)}{\sqrt\beta}\right].
 \label{eq:corrected-T}
\end{equation}
The stress can be written without the canceling Gaussian terms, since
$C_1-C_2+C_3=0$. Define
\begin{align}
 P_1(x,t)&=\frac12\left[
 e^{t+x}\operatorname{erfc}(\xi+\sqrt t)
 +e^{t-x}\operatorname{erfc}(\sqrt t-\xi)\right],\\
 P_a(x,t)&=\frac a2\left[
 e^{\gamma t-a x}\operatorname{erfc}(\xi-a\sqrt t)
 -e^{\gamma t+a x}\operatorname{erfc}(\xi+a\sqrt t)\right],\\
 R_\beta(t)&=r_\beta e^{-\beta t}+r_+e^{at}+r_-e^{-at},\\
 D_\beta(z)&=C_1 e^z\operatorname{erfc}(\sqrt z)
 +aC_2e^{\gamma z}\operatorname{erf}(a\sqrt z)
 +C_3\sqrt\beta e^{-\beta z}\operatorname{erfi}(\sqrt{\beta z})
 +R_\beta(z).
\end{align}
Then
\begin{equation}
 \mathcal S_\beta(x,t)=C_1P_1(x,t)+C_2P_a(x,t)
 -C_3\sqrt\beta J_\beta(x,t)+e^{-x}R_\beta(t)
 -H(t-\tau)D_\beta(t-\tau).
 \label{eq:corrected-stress}
\end{equation}
The delayed term is evaluated only for $t>\tau$ and extended continuously
at the front. The formulas are interpreted by limits at coincident poles,
including $\gamma=1$ and $\gamma=\beta^2$. The supplied numerical
implementation does not implement these general limiting cases; the
benchmark parameters in Table~\ref{tab:problem} avoid them. The unit-step stress is
$\mathcal S_0=\lim_{\beta\to0^+}\mathcal S_\beta$.
For the temperature this limit is explicitly
\begin{align}
 \mathcal T_0(x,t)={}&-e^{-x}
 +\frac{e^{t-x}}2\operatorname{erfc}(\sqrt t-\xi)
 +\frac{e^{t+x}}2\operatorname{erfc}(\sqrt t+\xi)\notag\\
 &+2\sqrt{\frac t\pi}e^{-\xi^2}
 -x\operatorname{erfc}(\xi).
\end{align}
All responses vanish for nonpositive time. For either field
$\mathcal F\in\{\mathcal T,\mathcal S\}$, the benchmark pulses follow by
linearity:
\begin{align}
 F_{\rm exp}(x,t)&=A_e\mathcal F_\beta(x,t),\\
 F_{\rm step}(x,t)&=A_s\left[\mathcal F_0(x,t)
 -H(t-\Delta)\mathcal F_0(x,t-\Delta)\right],\\
 F_{\rm smooth}(x,t)&=A\left[\mathcal F_{\beta_1}(x,t)
 -\mathcal F_{\beta_2}(x,t)\right].
\end{align}
For pulses nonzero at startup, the delayed stress response has a derivative
discontinuity at $x=at$. The rectangular pulse also launches a delayed
front from its switch-off time. The smooth pulse cancels the leading
startup discontinuity because $f(0^+)=0$.

For numerical evaluation, the identity
$e^{c^2t\pm cx}\operatorname{erfc}(\xi\pm c\sqrt t)
=e^{-\xi^2}\operatorname{erfcx}(\xi\pm c\sqrt t)$
is useful, but negative large arguments and cancellation between delayed
terms still require analytic regrouping. In particular, replacing every
factor individually by a scaled complementary error function does not
by itself remove all large intermediate terms. The released reference
implementation regroups the canceling $e^{\gamma t}$
contributions behind the front.

\FloatBarrier
\bibliographystyle{plainnat}
\bibliography{refs}
%TC:endignore
\end{document}
